% Options for packages loaded elsewhere
\PassOptionsToPackage{unicode}{hyperref}
\PassOptionsToPackage{hyphens}{url}
\documentclass[
]{article}
\usepackage{xcolor}
\usepackage{amsmath,amssymb}
\setcounter{secnumdepth}{-\maxdimen} % remove section numbering
\usepackage{iftex}
\ifPDFTeX
  \usepackage[T1]{fontenc}
  \usepackage[utf8]{inputenc}
  \usepackage{textcomp} % provide euro and other symbols
\else % if luatex or xetex
  \usepackage{unicode-math} % this also loads fontspec
  \defaultfontfeatures{Scale=MatchLowercase}
  \defaultfontfeatures[\rmfamily]{Ligatures=TeX,Scale=1}
\fi
\usepackage{lmodern}
\ifPDFTeX\else
  % xetex/luatex font selection
\fi
% Use upquote if available, for straight quotes in verbatim environments
\IfFileExists{upquote.sty}{\usepackage{upquote}}{}
\IfFileExists{microtype.sty}{% use microtype if available
  \usepackage[]{microtype}
  \UseMicrotypeSet[protrusion]{basicmath} % disable protrusion for tt fonts
}{}
\makeatletter
\@ifundefined{KOMAClassName}{% if non-KOMA class
  \IfFileExists{parskip.sty}{%
    \usepackage{parskip}
  }{% else
    \setlength{\parindent}{0pt}
    \setlength{\parskip}{6pt plus 2pt minus 1pt}}
}{% if KOMA class
  \KOMAoptions{parskip=half}}
\makeatother
\setlength{\emergencystretch}{3em} % prevent overfull lines
\providecommand{\tightlist}{%
  \setlength{\itemsep}{0pt}\setlength{\parskip}{0pt}}
\usepackage{bookmark}
\IfFileExists{xurl.sty}{\usepackage{xurl}}{} % add URL line breaks if available
\urlstyle{same}
\hypersetup{
  pdftitle={CSI interview year 1},
  hidelinks,
  pdfcreator={LaTeX via pandoc}}

\title{CSI interview year 1}
\author{}
\date{}

\begin{document}
\maketitle

\section{Introduction}\label{introduction}

\begin{itemize}
\tightlist
\item
  Context:

  \begin{itemize}
  \tightlist
  \item
    This thesis, forms part of a collaboration between CETIM and Inria.
    This PhD project aims to explore and develop model hybridization
    approaches applied to digital twins in mechanical engineering, by
    combining methods derived from physical modelling with data-driven
    approaches.
  \end{itemize}
\item
  Problem:

  \begin{itemize}
  \tightlist
  \item
    The Digital twins are build more and more and to achieve some
    qualify, make optimization, engineer combine, knowledge-based models
    (e.g., first-principles, physic based) called deductive and
    data-driven models (e.g., machine learning, statistical models)
    called inductive. But this approch is done ad-hoc and
    it\textquotesingle s.
  \item
  \end{itemize}
\end{itemize}

\section{Overview of the PhD project}\label{overview-of-the-phd-project}

\begin{itemize}
\tightlist
\item
  Research questions and hypotheses

  \begin{itemize}
  \tightlist
  \item
    What scientific community define hybridization, how the use it in
    their engineering process
  \item
    Inspired by the functionnement of design pattern, is it possible and
    relevant to have such a pattern for the hybridization ?
  \item
    How to make the hybridization works systematic ?
  \end{itemize}
\item
  Potential expected result

  \begin{itemize}
  \tightlist
  \item
    propose a systematic descriptive framework for conception of
    hybridization, and to identify the corresponding pattern to
    appplied, or to describe the engineering conception
  \item
    purprose a tools or DSLs for assisting engineer in the hybridization
    process, for pattern for example
  \end{itemize}
\end{itemize}

\section{Ongoing work}\label{ongoing-work}

\begin{itemize}
\item
  Literature review

  \begin{itemize}
  \tightlist
  \item
    Gang of 4: Design pattern

    \begin{itemize}
    \tightlist
    \item
      TRhe design pattern used in software engineering
    \end{itemize}
  \item
    Globalization Deamin-Specific lamuages

    \begin{itemize}
    \tightlist
    \item
      Show the importance of SL what to consider if relevant to create
      one
    \end{itemize}
  \item
    Designing an ISO 23247-compliant Hybrid Digital Twin Architecture
    for industry

    \begin{itemize}
    \tightlist
    \item
      Exposed an approch to consider Hybridization in the constructioin
      of Digital Twin (DT), as complement of the ISO 23247 which is one
      of standard for the realization of DT
    \end{itemize}
  \item
    Hybrid modeling design patterns:

    \begin{itemize}
    \tightlist
    \item
      Expose some pattern that are used frequently for hybriding models
    \end{itemize}
  \item
    Examining model qualities and their impact on digital twins:

    \begin{itemize}
    \tightlist
    \item
      Formaliize from mathematical point of view models properties and
      how their serve or contribute to the DT quality
    \end{itemize}
  \item
    GREY BOX MODELLING -- BRANCHES AND EXPERIENCES

    \begin{itemize}
    \tightlist
    \item
      Explore some case of combinaison, more specially some kind usage
      where both inductive and deductive models are used
    \end{itemize}
  \item
    Combining Machine Learning and Simulation to a Hybrid Modelling
    Approach: Current and Future Directions:

    \begin{itemize}
    \tightlist
    \item
      inductive and deductive models level of combinaion. First expose
      some main part of each model and show that each model can be
      combine with the other at different level exposed and the result
      is different
    \end{itemize}
  \end{itemize}
\item
  Timelines \& main results (see restitution format)\\
  Including:

  \begin{itemize}
  \tightlist
  \item
    Hybridization Characteristics (two dimensions identiifed)

    \begin{itemize}
    \tightlist
    \item
      Motivation: express why hybridizationn is done
    \item
      Realisatiomn: express how this hybrididization is implemented
    \end{itemize}
  \item
    Examplar paper ongoing - application of hybridization
    characteristics
  \end{itemize}
\end{itemize}

\section{Training hours (47 hours):}\label{training-hours-47-hours}

\begin{itemize}
\item
  Scientifics (49 hours)

  \begin{itemize}
  \tightlist
  \item
    Introduction à l'éthique et intégrité scientifique (3 hours)
  \item
    Master Class Science Ouverte (6 hours)
  \item
    Minor Quantum engineering (24) hours)
  \item
    SS01 : PSC1 Initial (7 hours)
  \item
    Worksops: MOMI (9 hours)
  \end{itemize}
\item
  Professionals (2 hours)

  \begin{itemize}
  \tightlist
  \item
    virtuous tryptic towards reproducibility (2 hours)
  \end{itemize}
\end{itemize}

\section{Next steps~~}\label{next-steps}

\begin{itemize}
\tightlist
\item
  Short term goal
\item
  Mid term goal
\item
  Long term goal
\end{itemize}

\section{Conclusion}\label{conclusion}

\end{document}
